\documentclass[11pt]{article}
\usepackage[margin=1in]{geometry}
\usepackage{amsmath,amssymb,amsthm}
\usepackage{graphicx}
\usepackage[hidelinks]{hyperref}
\setlength{\emergencystretch}{2em}

\newtheorem{theorem}{Theorem}
\newtheorem{lemma}{Lemma}
\newtheorem{corollary}{Corollary}

\newcommand{\Kcal}{\mathcal K}
\newcommand{\Prim}{\operatorname{Prim}}

\title{Finite-Dimensional Extensions of Nonseparable Compact-Operator Algebras Fail BCP}
\author{Run \texttt{fa\_banach\_001}, lane 19}
\date{June 17, 2026}

\begin{document}
\maketitle

\section*{Source Problem}

The source paper is Jinghao Huang, Karimbergen Kudaybergenov, and Rui Liu,
\emph{The ball-covering property of von Neumann algebras and noncommutative
symmetric spaces}, arXiv:2311.14257.  After proving that a commutative unital
\(C^*\)-algebra \(C(K)\) has BCP/UBCP exactly when its pure state space, i.e.
\(K\), has a countable \(\pi\)-basis, the authors ask for a noncommutative
version:

\begin{center}
\includegraphics[width=0.92\textwidth]{figures/open_problem_crop.png}
\end{center}

\section*{Statement}

The continuous-trace positive packet proved the expected criterion under a
separable-fibre hypothesis.  The following elementary subcase shows that this
hypothesis is necessary, even for a one-point primitive spectrum.  The same
argument also survives arbitrary finite-dimensional extensions of the form
\(\Kcal(H)+F\).

\begin{theorem}
\label{thm:main}
Let \(H\) be a nonzero Hilbert space, let \(F\subseteq\mathcal B(H)\) be a
finite-dimensional \(C^*\)-subalgebra, and put
\[
        A_F=\Kcal(H)+F.
\]
Then \(A_F\) has BCP if and only if \(H\) is separable.  Equivalently, \(A_F\)
has UBCP if and only if \(H\) is separable.
\end{theorem}

\begin{corollary}
\label{cor:type-i}
There is a type-I continuous-trace \(C^*\)-algebra with one-point primitive
ideal space and without BCP: take \(A=\Kcal(H)\) for any nonseparable Hilbert
space \(H\).  Thus countability of \(\Prim(A)\), or even the one-point
continuous-trace case, does not imply BCP unless the fibre is separable.

More generally, every finite-dimensional extension \(\Kcal(H)+F\) with
nonseparable \(H\) fails BCP.  Taking \(F=\mathbb C1\) gives a unital type-I
\(C^*\)-algebra with finite primitive ideal space and without BCP.
\end{corollary}

\begin{proof}
The algebra \(\Kcal(H)\) is elementary, hence type I and continuous trace.
It is simple for every nonzero Hilbert space \(H\), so \(\Prim(\Kcal(H))\) is
a singleton.  Each \(\Kcal(H)+F\) is an extension of a finite-dimensional
\(C^*\)-algebra by \(\Kcal(H)\), hence is type I and has finite primitive
ideal space.  The failure of BCP in the nonseparable case is
Theorem~\ref{thm:main}.
\end{proof}

\section*{Idea of the Proof}

The positive direction is just separability: if \(H\) is separable, then
\(\Kcal(H)\) is a separable Banach space, hence has UBCP.

For the negative direction, fix countably many proposed centers
\((T_n)\subset \Kcal(H)\).  Each compact operator has separable range, and so
does its adjoint.  The closed span of all ranges of \(T_n\) and \(T_n^*\) is
therefore separable.  In a nonseparable Hilbert space choose a unit vector
\(\xi\) orthogonal to that span.  The rank-one projection \(p_\xi\) is
orthogonal to every \(T_n\) on both sides:
\[
        p_\xi T_n=T_n p_\xi=0.
\]
Thus \(p_\xi-T_n\) is block diagonal with blocks \(1\) and \(-T_n\), so
\[
        \|p_\xi-T_n\|=\max(1,\|T_n\|)\ge \|T_n\|.
\]
No countable family of balls \(B(T_n,\|T_n\|)\) can cover the unit sphere.

For a finite-dimensional extension \(\Kcal(H)+F\), decompose the
finite-dimensional representation of \(F\) into matrix blocks with
multiplicity spaces.  Since \(H\) is nonseparable, one multiplicity space is
nonseparable.  Choose a vector \(\eta\) in that multiplicity space orthogonal
to the coordinate supports of the compact parts of the centers.  The
finite-dimensional slice over \(\eta\) reduces all compact parts and carries
the same \(F\)-matrix block as infinitely many nearby slices in the complement.
Thus the norm of each center is already visible on the complement of the
chosen slice, while a rank-one projection inside the chosen slice escapes the
corresponding BCP ball.

\section*{Proof}

We use the standard equivalent form of BCP: a normed space \(X\) has BCP iff
there are countably many centers \(x_n\in X\) such that
\[
        S_X\subseteq \bigcup_{n=1}^\infty B(x_n,\|x_n\|).
\]
Indeed, enlarging any open ball not containing the origin to radius equal to
the norm of its center still gives an open ball missing the origin.

\begin{lemma}
\label{lem:separable-range}
If \(T\in \Kcal(H)\), then \(\overline{T(H)}\) is separable.
\end{lemma}

\begin{proof}
The image \(T(B_H)\) of the unit ball is relatively compact.  Its closure is a
compact metric space, hence separable.  Since
\[
        T(H)=\bigcup_{m=1}^\infty m\,T(B_H),
\]
the closed linear span \(\overline{T(H)}\) is separable.
\end{proof}

\begin{lemma}
\label{lem:matrix-compact-lower}
Let \(E\) be a nonzero finite-dimensional Hilbert space and let \(K\) be an
infinite-dimensional Hilbert space.  Let \(Y\) be a Hilbert space containing
\(E\otimes K\) isometrically.  If \(S\in\mathcal B(E)\) and
\[
        R=(S\otimes1_K)+L:E\otimes K\to Y,
\]
where \(L:E\otimes K\to Y\) is compact, then
\[
        \|R\|\ge \|S\|.
\]
\end{lemma}

\begin{proof}
Choose a unit vector \(v\in E\) with \(\|Sv\|=\|S\|\), and choose an
orthonormal sequence \((\eta_m)\) in \(K\).  The sequence
\((v\otimes\eta_m)\) is weakly null, and compact operators send weakly null
sequences to norm-null sequences.  Hence
\[
        \|R(v\otimes\eta_m)\|
        \ge \|Sv\otimes\eta_m\|-\|L(v\otimes\eta_m)\|
        \to \|S\|,
\]
which gives the claimed lower bound.
\end{proof}

\begin{proof}[Proof of Theorem~\ref{thm:main}]
If \(H\) is separable, then \(\Kcal(H)\) is separable: finite-rank operators
with rational complex matrix entries relative to a fixed countable
orthonormal basis are norm dense in \(\Kcal(H)\).  Since \(F\) is
finite-dimensional, \(A_F=\Kcal(H)+F\) is separable.  Every separable normed
space has UBCP: choose a countable dense set \((d_j)\) in the unit sphere and
use centers \(2d_j\) with radius \(3/2\).  Hence \(A_F\) has UBCP, and
therefore BCP.

Assume now that \(H\) is nonseparable, and let \((a_n)\) be any countable
family of centers in \(A_F\).  Write
\[
        a_n=f_n+T_n,\qquad f_n\in F,\quad T_n\in\Kcal(H).
\]
By finite-dimensional representation theory, after a unitary identification
we may write
\[
        H=H_0\oplus\bigoplus_{r=1}^m (E_r\otimes K_r),
\]
where each \(E_r\) is finite dimensional, \(F\) acts as \(0\) on \(H_0\), and
on \(E_r\otimes K_r\) it acts as \(\rho_r(F)\otimes1_{K_r}\) for a
finite-dimensional representation \(\rho_r\).  If \(H_0\ne0\), include it as
the extra index \(r=0\) by putting \(E_0=\mathbb C\), \(K_0=H_0\), and
\(\rho_0=0\).  Since there are only finitely many summands and \(H\) is
nonseparable, choose an index \(r\) for which \(K_r\) is nonseparable.  Write
\[
        E=E_r,\qquad K=K_r,\qquad \rho=\rho_r.
\]

By Lemma~\ref{lem:separable-range}, the closed subspace
\[
        M=\overline{\operatorname{span}}\bigcup_{n=1}^\infty
        \bigl(T_n(H)\cup T_n^*(H)\bigr)
\]
is separable.  Let \(P_r\) be the projection onto the chosen summand
\(E\otimes K\).  Since \(E\) is finite dimensional, there is a separable
subspace \(N\subseteq K\) such that
\[
        P_rM\subseteq E\otimes N.
\]
Choose a unit vector \(\eta\in K\ominus N\), and put
\[
        V=E\otimes\mathbb C\eta.
\]
Then \(V\perp M\).  Hence \(T_nV=T_n^*V=0\), so \(V\) reduces every \(T_n\).
It also reduces \(F\), and the \(F\)-block on \(V\) is \(\rho(f_n)\).
Choose any rank-one projection \(p\) whose range is contained in \(V\).  Then
\(p\in\Kcal(H)\subseteq A_F\) and \(\|p\|=1\).

Let \(b_n=a_n|_{V^\perp}\).  Because \(a_n\) is block diagonal with respect to
\(H=V\oplus V^\perp\),
\[
        \|a_n\|=\max\bigl(\|\rho(f_n)\|,\|b_n\|\bigr).
\]
The space \(K'=K\ominus(N+\mathbb C\eta)\) is infinite dimensional.  On
\(E\otimes K'\subseteq V^\perp\), the operator \(b_n\) agrees with
\[
        \rho(f_n)\otimes1_{K'}+T_n|_{E\otimes K'}.
\]
The second term is compact, so Lemma~\ref{lem:matrix-compact-lower} gives
\(\|b_n\|\ge\|\rho(f_n)\|\).  Consequently \(\|a_n\|=\|b_n\|\).

Finally \(p-a_n\) is also block diagonal for \(H=V\oplus V^\perp\), with
blocks \(p|_V-\rho(f_n)\) and \(-b_n\).  Thus
\[
        \|p-a_n\|
        =
        \max\bigl(\|p|_V-\rho(f_n)\|,\|b_n\|\bigr)
        \ge \|b_n\|
        =
        \|a_n\|.
\]
The unit-sphere element \(p\) is not in any open ball
\(B(a_n,\|a_n\|)\).  Since the countable family of centers was arbitrary,
\(A_F\) fails BCP.
\end{proof}

\section*{Scope and Limitations}

This theorem does not solve the arbitrary \(C^*\)-algebra classification, but
it supplies a sharp type-I obstruction.  The previous continuous-trace theorem
required separable compact-operator fibres.  The present result proves that
dropping fibre separability is impossible: even the elementary algebra
\(\Kcal(H)\) has one primitive ideal when \(H\) is nonseparable, yet fails BCP.
The finite-dimensional extension clause strengthens the obstruction: adding
any finite-dimensional algebra \(F\subseteq\mathcal B(H)\) does not repair
BCP.

This also gives a simpler type-I analogue of the earlier one-point primitive
obstruction from \(\mathrm{II}_1\) factors.  The obstruction here is not
diffuseness; it is the availability of a new rank-one projection orthogonal to
the separable support of any countable list of compact centers.

\section*{Novelty and Search Bounds}

The run indexes were searched for arXiv:2311.14257, compact operators,
unitization, finite-dimensional extension, nonseparable Hilbert space,
\(\Kcal(H)\), one-point primitive space, finite primitive space, BCP, and
UBCP.  The existing packets included the continuous-trace separable-fibre
positive theorem and the \(\mathrm{II}_1\)-factor one-point primitive
obstruction, but no compact-operator fibre-separability obstruction.  Web
searches on June 17, 2026 for BCP and compact operators found the source paper
and adjacent operator-space BCP literature, but no exact statement of
Theorem~\ref{thm:main}.

\section*{Human Review Focus}

The proof is low-tech.  Reviewers should check Lemma~\ref{lem:separable-range}
for compact operators on nonseparable Hilbert spaces,
Lemma~\ref{lem:matrix-compact-lower} for compact perturbations of a
finite-dimensional matrix block, the finite-dimensional representation
decomposition of \(F\), and the final block-diagonal norm computation.  Once
those are accepted, the BCP failure is immediate.

\begin{thebibliography}{9}

\bibitem{HKL}
Jinghao Huang, Karimbergen Kudaybergenov, and Rui Liu,
\emph{The ball-covering property of von Neumann algebras and noncommutative
symmetric spaces},
arXiv:2311.14257, 2023.

\bibitem{lane19separable}
Run \texttt{fa\_banach\_001}, lane 19,
\emph{The Separable \(C^*\)-Algebra Case of BCP},
solution packet
\texttt{2311.14257\_separable\_cstar\_bcp\_classification}, 2026.

\bibitem{lane19ct}
Run \texttt{fa\_banach\_001}, lane 19,
\emph{BCP for Continuous-Trace \(C^*\)-Algebras with Separable Fibres},
solution packet
\texttt{2311.14257\_continuous\_trace\_separable\_fiber\_bcp}, 2026.

\end{thebibliography}

\end{document}
